\clearpage\phantomsection\addcontentsline{toc}{section}{校勘说明（译者附记）}\section*{校勘说明（译者附记）}

本附记为中译重建时加入的校勘材料，不属于原书正文。十篇已出版报告、卷首及索引均由未负责相应初译的独立代理检查源译覆盖、假设和证明步骤；所列疑点又逐项核对了法文原页。这里区分英文重构错误、法文原书错误、成立但需要解释的读法，以及证据不能唯一复原的原文疑义。可确定的修订已用于正文，原读法与依据记在各条中。原版与英文本均未改动。

本轮最初记录 64 项，复核新增 18 项，共 82 项均有裁定；没有未处理的清单项。仍有明确保留范围的原文疑义为 I-01、III B-20、XII 的(6.4.7)及导言的概述编号，详见对应条目。正文上标“[校]”可跳至相关说明。独立代理检查不构成专业人工审定；本次也不宣称逐句核对了整部法文本或重新证明了全部定理。以下“原书页码”均指法文印刷页，必要时另列物理PDF页。

\hypertarget{collation-I-01}{}

\paragraph*{I-01：定理 2.1 证明的失效引文。}法文印本第17页（PDF第28页）实作“en vertu de 1.9 b)(ii)”，随后对 $F\in\ob\D_c(X)$ 写出 $F\otimes^{\mathrm L}P\simeq\underline D'\underline DF$。同书第10页的1.9只有(a)、(b)，其中(b)是逆像保持可构造性的引理，并无(ii)；英文重构本额外称其为“Proposition”。命题1.11包含单一对偶函子的张量与内Hom交换公式，但不能仅凭编号相近，就把这里的双对偶公式及所写的整个 $\D_c(X)$ 范围直接归给其中某款。因此保留原印本引文，明确记为失效引用；现有印本证据不能唯一恢复原拟指的条款。本条未据此宣称公式(**)或定理错误，也未宣称已经补足其在无界范围上的证明。

\hypertarget{collation-I-02}{}

\paragraph*{I-02：对偶化定义的编号。}法文印本第32页（PDF第43页）在定理3.4.1证明首句实作“conditions (i) à (iii) de (1.6)”。第8页（PDF第19页）的定义1.7才列出这三个条件；第7页的命题1.6只论拟有限非常逆像的拟内射维数。正文据此将该处1.6校为1.7，未改变三个条件或定理假设。

\hypertarget{collation-I-03}{}

\paragraph*{I-03：正像保持双对偶的三处引用。}法文印本第32、33、34页（PDF第43、44、45页）分别印作“En vertu de (1.11)”、“Par (1.11)”和“Appliquant à nouveau (1.11)”。第14页（PDF第25页）的推论1.13恰好断言：固有态射的导出正像保持双对偶；有限态射还可反推双对偶。这里 $g=fh$ 为固有态射，另两处闭浸入均有限，故正文三处均校为推论1.13。命题1.11论张量与内Hom的交换，不是这三步所用的正像断言。

\hypertarget{collation-I-04}{}

\paragraph*{I-04：系数环与指数。}法文印本第45页（PDF第56页）可辨的原读法是 $A=\Z/\ell\Z$ 与“pour $\nu\geqslant2$”；环式中没有实际印出指数 $\nu$。英文重构本把条件也误读成了 $\ell\geqslant2$。本报告第16页（PDF第27页）曾明确使用素数幂记号 $A=\Z/\ell^\nu\Z$、$\ell$ 为素数，但该页的局部环假设仅限第2节，不能直接冒称第4节已有这一假设。正文把本句限定为素数幂系数的情形，并校读作 $A=\Z/\ell^\nu\Z$、$\nu\geqslant2$：对这种环，$\Z/\ell\Z$ 有以乘 $\ell$ 与乘 $\ell^{\nu-1}$ 交替为微分的无限自由预解，故各正次数的 $\operatorname{Ext}_A^i(\Z/\ell\Z,\Z/\ell\Z)$ 均不为零；反之，$\nu=1$ 时 $A$ 为域，全局维数为零。这是依据上下文及代数事实作出的明确校订，缺失的指数不是扫描本已见字形。1977年增补纠正的是导出张量积可定义的范围，并未纠正本行环式；该增补仍完整保留。

\hypertarget{collation-I-05}{}

\paragraph*{I-05：中间同构的移位。}法文印本第46页（PDF第57页）确实印作 $\R\Hom(A_{V,X},F'_{U,X})[1]\simeq\R\Gamma_x(F'_{U,X})$；第47页的(4.7.10)却使用 $[-1]$。令 $E=F'_{U,X}$，由 $A_{V,X}\to A_X\to A_x\to A_{V,X}[1]$ 作用反变函子 $\R\Hom(-,E)$，并用 $\R\Hom(A_X,E)=\R\Gamma(X,E)=0$，得 $\R\Hom(A_{V,X},E)\simeq\R\Hom(A_x,E)[1]$。故所需式必为 $\R\Hom(A_{V,X},E)[-1]\simeq\R\Gamma_x(E)$，与(4.7.10)一致；正文仅将中间式的 $[1]$ 校为 $[-1]$。

\hypertarget{collation-I-06}{}

\paragraph*{I-06：推论 1.13 的原同构链成立。}法文印本第14页（PDF第25页）的中间项确为 $\underline D_Y\R f_*\underline D_XF$，两个箭头确实都引用(a)(i)。由于推论假设 $f$ 固有，$\R f_*=\R_!f$，故两次使用命题1.12(a)(i)即可得到原链，无需改动任何正像符号。英文重构本另加的编者注曾建议中间改为 $\R_!f$、末箭头改引(a)(ii)；前者在此与原式相同，后者不必要，且(a)(ii)要求两个 $K$ 都为对偶化复形，这并非本推论假设。正文保留原链；原英文编者注的建议及本次否定理由移记于此，不把正确的法文原证明标为错误。

英文重构稿原编注的完整译文（其判断已在上文复核）：
\begin{quote}\small
\emph{编者注.} 印本原文的中间项为 $\underline D_Y\R f_*\underline D_XF$, 并对相继的两个箭头均引用命题~1.12(a)(i). 本版保留该同构链. 数学上一致的修订是在中间一行写 $\underline D_Y\R_!f\,\underline D_XF$, 并对最后一个箭头引用命题~1.12(a)(ii).
\end{quote}

\hypertarget{collation-I-07}{}

\paragraph*{I-07：4.6 普通诱导行的读法与对象所在空间。}法文印本第43页（PDF第54页）经放大可辨，第三行第一个拉回号为 $i'^*$，对偶函子下标实作大写 $X$，即 $i'^*\underline D_Xi''^*\R^!j(F)$；它与同页下一行的 $\underline D_{\bar x}$ 字形不同。英文重构本及其旧编者注把此处保留为 $\underline D_{\bar x}$，又称首个拉回号不清楚，本次图像证据不支持这两项读法。令 $C=\overline{\{x\}}$；由于 $i'':C\to Y$，$i''^*\R^!j(F)$ 是 $C$ 上的对象，普通诱导公式1.12(b)(i)的中间项必须使用 $\underline D_C$，随后才经 $i'^*$ 和4.3变为 $\underline D_{\bar x}$。正文据对象所在空间校为 $i'^*\underline D_{\overline{\{x\}}}i''^*\R^!j(F)$；此下标是依据公式类型的校订，原印本大写 $X$ 仍在本条记录。

英文重构稿原编注的完整译文（其判断已在上文复核）：
\begin{quote}\small
\emph{编者注.} 印本原文中通常诱导公式所在的一行无法无歧义地辨认. 当前经原文核对的法文 TeX 写为 $i'^*\underline D_{\ol{x}}i''^*\R^!j(F)$; 但沿 $i'':\ol{\{x\}}\to Y$ 检查对象所属的范畴, 下标应当是 $\ol{\{x\}}$, 而原文字形又可能看成 $X$, 第一个拉回记号也不清楚. 本版保留法文校核本的读法, 并不声称扫描本足以作出裁决.
\end{quote}

\hypertarget{collation-I-08}{}

\paragraph*{I-08：引理 2.2 中相反的移位符号。}法文印本第19页（PDF第30页）先令 $L=\uH^{-r_0}(P)$、$L'=\uH^{r_0}(P')$，却随后印作 $P\simeq L[-r_0]$、$P'\simeq L^\vee[+r_0]$。同书第18页还明确约定 $P_{\bar x}\simeq A_{\bar x}[r(x)]$。由 $\uH^i(C[r])=\uH^{i+r}(C)$，且上述两复形各只剩一个非零上同调，必为 $P\simeq L[r_0]$、$P'\simeq L^\vee[-r_0]$。正文校正这两个符号，保留 $r_0$ 和 $L$ 的定义。

\hypertarget{collation-I-09}{}

\paragraph*{I-09：曲线证明中补入的平凡作用因子。}法文印本第51页（PDF第62页）在(5.1.4)之后也印作“le facteur $\widehat{\Z}(1)$”，并非只有英文丢了下标。由同节定义 $H=\prod_{\ell\ne\pi}\Z_\ell(1)$，这里实际把 $N$ 的 $H$-作用沿自然投射 $\widehat{\Z}(1)\to H$ 拉回；当 $\pi>0$ 时，其额外的核因子是 $\Z_\pi(1)$，应令这个因子平凡作用，而非令整个 $\widehat{\Z}(1)$ 平凡作用。因 $n$ 与 $\pi$ 互素，核的高次上同调为零，Hochschild--Serre 谱序列遂给出所述同构；$\pi=0$ 时没有额外因子。正文据此明写投射和核因子，原印本记号保留于本条。

\hypertarget{collation-I-10}{}

\paragraph*{I-10：3.3.1 证明中所验证条件的编号。}法文印本第29页（PDF第40页）在“(ii bis) $\Rightarrow$ (i bis)”段实作“vérifier (ii) pour $F=i_!(M)$”。此段正在由所有浸入的正像可构造性证明 $\R\underline{\Hom}(F,G)$ 可构造，接下来也确实计算 $\R\underline{\Hom}(i_!M,G)$；因而所验证的应为(i bis)，不是已经由假设包含的(ii)。正文仅将这一处(ii)校为(i bis)。

\hypertarget{collation-I-11}{}

\paragraph*{I-11：纯性计算中 Tate 扭转的符号。}法文印本第27页（PDF第38页）的注3.2.2实作 $\underline{\operatorname{Ext}}^{2d}(A_U,A_X)|_U=(\mu_n)_U^{\otimes d}$，第28页（PDF第39页）的推论3.2.5证明也实作 $(\mu_n)^{\otimes d}(k)$。但同书第23页（PDF第34页）定义3.1.4的局部基本类同构为 $A_Y\simeq\underline{\mathrm H}^{2d}_Y((\mu_n)_X^{\otimes d})$；消去这一可逆扭转，就得到 $\underline{\mathrm H}^{2d}_Y(A_X)\simeq(\mu_n)_Y^{\otimes(-d)}$。因此这两处由纯性得到的自然同构，其右端应分别为 $(\mu_n)_U^{\otimes(-d)}$ 与 $(\mu_n)^{\otimes(-d)}(k)$；正文补入负号。第二处在可分闭域上正负扭转虽非典范地同构，故原来的非零性结论不受影响，但按所引纯性同构应使用负扭转。原印本两处正号在此保留。

\hypertarget{collation-III-01}{}

\paragraph*{III-01。}

此处次数平移据迹箭头及(3.2.4)改为 $[-2N]$; 法文印本第90页文字误作 $[2N]$.

\hypertarget{collation-III-02}{}

\paragraph*{III-02。}

此处拉回记号据法文印本第95页恢复为 $f_2^*$; 英文重构稿误作 $f_2^!$.

\hypertarget{collation-III-03}{}

\paragraph*{III-03。}

这里的推论编号据上下文改为3.5; 法文印本第124页误作6.5.

\hypertarget{collation-III-04}{}

\paragraph*{III-04。}

此处范畴下标据 $f_i:X_i\to Y_i$ 及对象的定义域改为 $Y_i$; 法文印本第104页误作 $X_i$.

\hypertarget{collation-III-05}{}

\paragraph*{III-05。}

法文印本第121页这里的结构层下标误作 $X,Y$; 按正在讨论的两个概形改为 $X_1,X_2$.

\hypertarget{collation-IIIB-01}{}

\paragraph*{IIIB-01。}

法文印本第143页这里的目标误作 $(DM)_x$; 按对偶对应的定义改为 $(DL)_x$.

\hypertarget{collation-IIIB-02}{}

\paragraph*{IIIB-02。}

法文印本第144页此处误作 $\{y\}\times Y$; 按 $x\in X,y\in Y$ 及1.2改为 $\{x\}\times Y$.

\hypertarget{collation-IIIB-03}{}

\paragraph*{IIIB-03。}

法文印本第151页 $f$ 的定义域误作 $T-Z$; 由上行的基变换及下行的复合改为 $T-Y$.

\hypertarget{collation-IIIB-04}{}

\paragraph*{IIIB-04。}

引理4.3的包含图及其后有关 $j^\prime,m^\prime$ 的公式均据法文印本第158--159页恢复: $j^\prime:Z-(X\cup Y)\to Z-X$, $m^\prime:Y-z\to Z-X$. 英文重构稿漏去撇号、反置下方两箭头, 并将若干 $Rj^\prime_*L$ 误写为 $Rj_*i^\prime_!L$.

\hypertarget{collation-IIIB-05}{}

\paragraph*{IIIB-05。}

法文印本第160页在证明c)时沿用了a)中的 $\bar z$; 由c)中 $\widetilde C$ 的定义及右下方 $\bar s$ 处的茎, 此处应在 $\bar s$ 处严格局部化.

\hypertarget{collation-IIIB-06}{}

\paragraph*{IIIB-06。}

法文印本第161页(*)右端的空间及限制下标有排印错误. 这里按原命题的目标 $B-t$ 在爆破后的逆像 $B^\prime-D$ 修订; 后文各 $B^\prime_j-b_j$ 的不交分解也给出同一目标.

\hypertarget{collation-IIIB-I-NEW-01}{}

\paragraph*{IIIB-I-NEW-01。}

法文印本第159页此处印作 $R^qj^\prime_*L=0$「pour $p>2$」. 由(P)及本段的分步归约, 应为 $q\ge2$; 这也与紧接着只处理 $q\le1$ 相符.

\hypertarget{collation-IIIB-07}{}

\paragraph*{IIIB-07。}

编校注：法文原版第167页以等号引入整个迹箭头；这不改变其定义域、值域或构造。若将等号规范为冒号，仅属记号整理。

\hypertarget{collation-IIIB-08}{}

\paragraph*{IIIB-08。}

编校注：按法文原版第171页恢复：有界性条件中的范畴为 $D^b(B^\circ)$；图左列上、下两项的张量积分别取在 $A$、$B$ 上。英文重构稿误加了包络代数的上标。

\hypertarget{collation-IIIB-09}{}

\paragraph*{IIIB-09。}

编校注：法文原版第172页此处作“竖直箭头”，但紧随的图中 $\Tr_A$、$\Tr_B$ 均为水平箭头，故据图改正。

\hypertarget{collation-IIIB-10}{}

\paragraph*{IIIB-10。}

编校注：法文原版第173页此处将乘积 $fg$ 误排为 $f_g$；据所引循环迹公式5.8.5改正。

\hypertarget{collation-IIIB-11}{}

\paragraph*{IIIB-11。}

编校注：按法文原版第184页恢复6.2.1左端的两个普通正像符号 $f_{1*}$、$f_{2*}$；右端原为 $f_*$。英文重构稿将左端两处星号误作叹号。

\hypertarget{collation-IIIB-12}{}

\paragraph*{IIIB-12。}

编校注：法文原版第184页将第二对象的右系数误写为 $B$；据上一行张量函子的定义域及结果的右系数 $C$，改为 $C$。

\hypertarget{collation-IIIB-13}{}

\paragraph*{IIIB-13。}

编校注：法文原版第185页此处作 $E_2=KA_{X_1}$；6.2.3中相应对象为 $F_1$，故改正对象名称。

\hypertarget{collation-IIIB-14}{}

\paragraph*{IIIB-14。}

编校注：按法文原版第187页恢复施行的函子 $i_*^ci^{c!}$，并删去6.5.7中 $c_*$ 上误加的上标 $c$；两处均为英文重构造成的记号损坏。

\hypertarget{collation-IIIB-15}{}

\paragraph*{IIIB-15。}

编校注：法文原版第194页此处的 $M_n$ 未定义；本段对象及紧随的迹均使用 $L_n$，故改为 $f_*L_n$。

\hypertarget{collation-IIIB-16}{}

\paragraph*{IIIB-16。}

编校注：法文原版第194页段首将开浸入误记为 $j$；按6.14.1及后续 $i_!\Lambda_U$ 的一致记号，改记为 $i$。

\hypertarget{collation-IIIB-17}{}

\paragraph*{IIIB-17。}

编校注：法文原版第195页此处作 $n\ge1$，与上一行 $F^1\Lambda_X=i_!\Lambda_U$ 及后面的 $\operatorname{gr}^1$ 项矛盾；改为 $n\ge2$。

\hypertarget{collation-IIIB-18}{}

\paragraph*{IIIB-18。}

编校注：法文原版第195页在 $\operatorname{gr}^1$ 行将目标误写为闭补集项；据该分次项为 $\Lambda_{U,X}$，改为 $c_2^!\Lambda_{U,X}$。

\hypertarget{collation-IIIB-19}{}

\paragraph*{IIIB-19。}

编校注：法文原版第195页左端第二下标误作 $U$；右端沿 $j:X-U\to X$ 取正像，故左端应为 $\underline c_{X-U,X}$。

\hypertarget{collation-IIIB-20}{}

\paragraph*{IIIB-20。}

编校注：法文原版第196页确作“$G$在 $C$ 和 $D$ 上平凡作用”。6.11已要求在 $Y,D$ 上平凡作用，而6.23又引用本段；将此处 $C$ 改为 $Y$ 虽可疑似消除衔接问题，却会实质减弱假设，尚无足够证据断定，故保留原文。

\hypertarget{collation-IIIB-21}{}

\paragraph*{IIIB-21。}

编校注：法文原版第197页此处作 $f_*M$；因 $M$ 已在 $Y$ 上、$L$ 在 $X$ 上，并据伴随箭头及6.18.3，改为 $f_*L$。

\hypertarget{collation-IIIB-22}{}

\paragraph*{IIIB-22。}

编校注：按法文原版第199页恢复完整交换图：补入 $h_*c_2^*L$、箭头 $h_*u$、编号(1)--(5)及方块(A)、(B)。英文重构稿漏掉了一个对象及相应箭头，不只是漏标编号。

\hypertarget{collation-IIIB-23}{}

\paragraph*{IIIB-23。}

编校注：法文原版第200页在同构 $\alpha$ 中将 $U$ 两次误写为 $X$；因 $f_V:U\to V$，该同构属于 $D(G,U)$，右端应为 $\Lambda_U\otimes^{\mathrm L}_{\Lambda}E_0$。

\hypertarget{collation-IIIB-24}{}

\paragraph*{IIIB-24。}

编校注：这里的归约还隐含对态射的重新记名：令 $H=Z_g$、$a=g f_*\underline c_{U,X}$、$b=gv_0$，则 $a,b$ 均为 $H$-等变态射，6.23.2成为恒等元坐标上的张量迹乘法公式；下述引理对任意这样的态射证明该式。仅把群换成中心化子本身并不会使元素 $g$ 等于 $e$。

\hypertarget{collation-IIIBII-NEW-01}{}

\paragraph*{IIIBII-NEW-01。}

编校注：按法文原版第184页恢复6.2.2中两个逆像函子的下标：分别为 $f_1^!M_1$、$f_2^!M_2$。

\hypertarget{collation-IIIBII-NEW-02}{}

\paragraph*{IIIBII-NEW-02。}

编校注：按法文原版第185页恢复6.3.1为 $t^!A=KA_T$；$A$ 是函子 $t^!$ 的输入对象，英文重构稿误将其放到下标。

\hypertarget{collation-IIIBII-NEW-03}{}

\paragraph*{IIIBII-NEW-03。}

编校注：法文原版第170页此处将值域写作 $A$；在本段的非交换情形，所引5.6.3及5.7.2的值域均为 $A_\natural$，故补上该下标。

\hypertarget{collation-IIIBII-NEW-04}{}

\paragraph*{IIIBII-NEW-04。}

编校注：法文原版第175页在5.10.8前的求值式左端误写 $\otimes^{\mathrm L}_{B^e}E$；按5.3.5及紧随的5.10.8，应为 $\otimes^{\mathrm L}_BE$。右端 $B\otimes^{\mathrm L}_{B^e}B$ 的下标仍保留。

\hypertarget{collation-IIIBII-NEW-05}{}

\paragraph*{IIIBII-NEW-05。}

编校注：法文原版第191页在可构造子范畴的等式左端漏写群变量 $G$；据上一等式及下一行的记号，补为 $D_c(G,{}_AX)$。

\hypertarget{collation-VVI-01}{}

\paragraph*{VVI-01。}

校注: 法文印刷页 213 在此处及下式均印作 $n\ge r-1$. 由 $Y=T[r]$ 及所述态射集的双射, 应在 $0\le n<r$ 时补零, 故两处改为 $n\le r-1$; 当 $n\ge r$ 时使用 $Y_{n-r}$ 及其转移态射.

\hypertarget{collation-VVI-02}{}

\paragraph*{VVI-02。}

校注: 法文印刷页 208 将复合写作 $w_{ts}\circ w_{sr}$. 根据上一页给出的源和目标及本报告的复合约定, 此处应为 $w_{sr}\circ w_{ts}$.

\hypertarget{collation-VVI-03}{}

\paragraph*{VVI-03。}

校注: 法文印刷页 215 在这两个平移后的复合中, 第一项遗漏了 $f$ 后的 $[r]$. 由平移的函子性及 $v_{rT}[r]=v_{rY}$ 补正.

\hypertarget{collation-VVI-04}{}

\paragraph*{VVI-04。}

校注: 法文印刷页 246 将下面两个公式及 $(d_1)$ 中的 $Z_r,B_r$ 名称互置. 按谱序列的循环与边界定义, 并与前文 b) 的稳定范围及下一页 $(d_2)$ 核对后, 此处同步校正四个名称.

\hypertarget{collation-VVI-05}{}

\paragraph*{VVI-05。}

校注: 法文印刷页 262 此处的指数为 $n+1$. 按本段对 $F_n$ 局部自由性的定义, $F_{n+1}$ 的系数环应为 $\mathbb Z/\ell^{n+2}\mathbb Z$, 故校正为 $n+2$.

\hypertarget{collation-VVI-06}{}

\paragraph*{VVI-06。}

校注: 法文印刷页 268 此处写作“仍局部常值”. 依本报告 1.2.5 的注, 局部常值与扭常值须加区分; 张量表示保证的是仍为扭常值层, 并不保证单值群的像有限. 此处按该报告自身的术语校正, 不增添任何假设.

\hypertarget{collation-VVI-07}{}

\paragraph*{VVI-07。}

校注: 法文印刷页 269 断言所选的 \(\mathbb Z_\ell\) 层有唯一的直和分解 \(F=P\oplus T\). 这一分解对底层模存在, 但一般不能选为与基本群作用相容的分解. 本处改用典范挠子层的无挠商, 即可得到同一个有理系数对象, 并完成下述同态层的构造.

校注: 法文印刷页 270 同样断言相应的整系数层可分裂为无挠部分与挠部分的直和. 如 1.4.6 的校注所说明, 不能一般地要求该分解与基本群作用相容; 应使用典范的无挠商.

例如令循环群 $C_\ell$ 在 $M=\Z_\ell\oplus\F_\ell$ 上由 $g(a,t)=(a,t+\bar a)$ 作用，则 $g^\ell=1$，挠子模为 $T=0\oplus\F_\ell$。若商映射有等变截面，$1$ 的像必为 $(1,c)$，但 $g(1,c)=(1,c+1)$，矛盾。这说明原书所需的等变直和分裂一般不存在；商 $M/T$ 的构造不需要这种分裂。

\hypertarget{collation-VVI-08}{}

\paragraph*{VVI-08。}

校注: 法文印刷页 270 用语为“有限型扩张”. 此处承接 1.4.3 的有限 \(\mathbb Q_\ell\) 系数代数语境, 所需的是有限代数扩张, 此时 \(A[1/\ell]=L\). 若允许 \(\mathbb Q_\ell(t)\) 等超越扩张, 后述商范畴不再给出 \(L\) 层范畴; 故在中译中明确有限扩张这一隐含限制.

\hypertarget{collation-VVI-09}{}

\paragraph*{VVI-09。}

校注: 法文印刷页 216 此处误作“同构的集合”. 按商范畴的态射集定义及随后函子 \(m\) 的构造, 应为全部“态射的集合”.

\hypertarget{collation-VVI-10}{}

\paragraph*{VVI-10。}

校注: 法文印刷页 267 在允许 \(L\) 为任意有限 \(\mathbb Q_\ell\) 代数之后, 仍使用“有限维 \(L\) 向量空间”的措辞. 一般应指底层为有限维 \(\mathbb Q_\ell\) 向量空间的有限型 \(L\) 模, 无须对 \(L\) 自由; 当 \(L\) 为域时才恢复通常的向量空间说法.

\hypertarget{collation-VVI-11}{}

\paragraph*{VVI-11。}

校注: 法文印刷页 251 此处引作 V 3.2.1. 实际的 \(J\) 进投射系定义为 V 3.1.1; V 3.2.1 是 AR-\(J\) 进性的判据, 故校正此引用.

\hypertarget{collation-VVI-12}{}

\paragraph*{VVI-12。}

校注: 法文印刷页 271 此处可见的引用为“5 1.3”; 英文重构作“V 1.3”. 按诺特 AR-\(J\) 进对象的定义, 应规范为 V 5.1.3.

\hypertarget{collation-VIIVIII-01}{}

\paragraph*{VIIVIII-01。}

VII 7.3 证明：原书第322页作“(3.10)”。第3节止于3.9；此处引用的是第304页推论4.9的公式 $\operatorname{EP}_\ell(Y)=p_{Y*}c_{\dim Y}(\check\Omega_Y^1)$，故改为“(4.9)”。

\hypertarget{collation-VIIVIII-02}{}

\paragraph*{VIIVIII-02。}

VIII 8.1 的证明标题：原书第367页误作“命题5.2的证明”；第366页的待证命题及证明结尾均为8.1，故改为8.1。

\hypertarget{collation-VIIVIII-03}{}

\paragraph*{VIIVIII-03。}

VII 2.2 构造 $c_i$：原书第287页作“(2.2.1)”。张量积后所用的乘法是第285--286页的(2.1.1)，而定理2.2.1尚在下文，故改为(2.1.1)。

\hypertarget{collation-XXII-01}{}

\paragraph*{XXII-01。}

X 3.6：原书第379页作“在4.1中定义”。带附加模结构的 $\RHom$ 构造在3.4；4.1讨论不动点重数，故改为3.4。

\hypertarget{collation-XXII-02}{}

\paragraph*{XXII-02。}

X 3.7：原书第380页作“(4.3.1)”。由(3.7.1)的 $g_{\check P}={}^t(g^{-1})_P$ 及转置保持迹得到(3.7.2)，故改为(3.7.1)。

\hypertarget{collation-XXII-03}{}

\paragraph*{XXII-03。}

X 3.9 证明：原书第382页作“(4.3.4)”，应为前页命题3.8的同构链(3.8.1)，据此改正。

\hypertarget{collation-XXII-04}{}

\paragraph*{XXII-04。}

X 4.2：原书第384页作“引理5.1”。素元差的赋值公式来自第382页引理4.1，故改为4.1。

\hypertarget{collation-XXII-05}{}

\paragraph*{XXII-05。}

X 定理7.1：原书第396页作“7.1的条件i)、ii)”。这两项条件列于第393--394页6.1，故改为“6.1的条件i)、ii)”。

\hypertarget{collation-XXII-06}{}

\paragraph*{XXII-06。}

XII (6.4.7)：原书第437页写作 $S'$，并把“有不动提升”时的值写为 $0$，“其他情形”写为 $\nu_x(f)-1$，与同页(6.4.8)及定义(4.1)冲突。这里改用前文已定义的 $S$，交换两值，并把提升限定在 $p^{-1}(x)$。该分段计算应在 $p$ 于 $x$ 上方平展的情形理解：有不动提升时恰有 $c^\varphi(g)$ 个，且每个重数为 $\nu_x(f)$；没有提升时定义中的和为空。原文仅假设 $F$ 在不动点不分歧，此条件本身不保证任意选定的平凡化覆盖 $p$ 在该处平展。因此，这个计算不能单独充当任意分歧覆盖下(6.4.6)的证明；原断言的一般性须另作论证。另见(6.3.2)的符号校正。

原文分段式完整保留如下，以便区分原读与上述限定范围后的校正：
\[
S'_{x,\ell^n}(f,f',G,\varphi)(g)=
\begin{cases}
0&\text{若存在 }x'\in X'\text{ 使 }(g^{-1}f')(x')=x'\text{, 或}\\
&\text{等价地, 若 $g$ 的 $\varphi$-共轭类等于}\\
&\text{类 }\mathrm{NW}_{f}^{G,\varphi}(x)\text{ (参见 3); }\\
\nu_x(f)-1&\text{其他情形.}
\end{cases}
\]

\hypertarget{collation-XXII-07}{}

\paragraph*{XXII-07。}

报告XI的引文沿用原书的历史编号。原书导言第VI页说明，I. Bucur 的稿件在搬家时遗失且无副本；III B 5（第163页）再次说明此事，并指向 SGA $4^{1/2}$ Rapport。现存更精细的理论见III B 5、6。本译本没有因此遗漏一篇可供翻译的报告。

\hypertarget{collation-XXII-08}{}

\paragraph*{XXII-08。}

XII (6.3.2)：原书第431页右端作 $\Tr_A(u_x)-\alpha_x$。由(6.3.4)、(6.3.5)、(6.3.11)、(6.3.12)及(6.2.2)的定义，右端应为 $\Tr_A(u_x)+\alpha_x$；这一符号也使不分歧情形得到(6.4.6)，故将减号改为加号。

\hypertarget{collation-XXII-09}{}

\paragraph*{XXII-09。}

X 2.2 证明中的 pseudo-cohérent 统一译为“伪凝聚”，以区别 quasi-cohérent 的“拟凝聚”。原书第375页两个位置均为 pseudo-cohérent。

\hypertarget{collation-XXII-10}{}

\paragraph*{XXII-10。}

X 2.2 证明后的预告：原书第375页作“Weil公式6.1”及“第5节定义的局部不变量”。现有正文中Weil公式为命题5.1（式5.2），其Swan局部表示定义在第4节，故相应改为5.1及第4节。

\hypertarget{collation-XV-01}{}

\paragraph*{XV-01。}

校勘注：法文原书第 472 页此处作 $F$，据前后逐层的记号补为 $F_\nu$。

\hypertarget{collation-XV-02}{}

\paragraph*{XV-02。}

校勘注：法文原书第 470 页此处的 $\bar F$ 漏下标 $\nu$；本句讨论逐层的复形，补为 $\bar F_\nu$。

\hypertarget{collation-XV-03}{}

\paragraph*{XV-03。}

校勘注：法文原书第 475 页括号内误作 $D^-(A)$；这里定义的是约化后的 $A_0$-模复形的伪凝聚性与完美性，应为 $D^-(A_0)$。

\hypertarget{collation-XV-04}{}

\paragraph*{XV-04。}

记号说明：本段原书写作 $\mathcal J^p$，这里指由 $\mathcal J$ 的局部截面的第 $p$ 次幂生成的理想，今常记为 $\mathcal J^{[p]}$；一般不同于通常的理想乘积幂。

\hypertarget{collation-XV-05}{}

\paragraph*{XV-05。}

校勘注：法文原书第 459 页这两个复合的因子次序颠倒；此处按态射的源与目标调正，亦与第 460 页证明中下方三角形所用的等式一致。

\hypertarget{collation-XV-06}{}

\paragraph*{XV-06。}

索引说明：原书的 XIV 与正文 XV 为同一报告（原目录作 XIV=XV）；索引中的第 20、21 页沿用原分篇页码，分别对应法文合订本第 461、462 页。下列现编定位按正文的节次列出；有理性定理原索引作 XIV 2 no.2，现定位为 XV §3 第2节。

\hypertarget{collation-FRONT-01}{}

\paragraph*{FRONT-01。}

校勘说明：法文导言第 V 页介绍导出范畴有限性与 Grothendieck 群的句子误作第 VII 报告，据原目录及第 VIII 报告正文改为 VIII。同页此前介绍经典概形、Chern 类和自交公式的 VII 正确。第 III 页总述对偶形式体系时原作 (I,VII)，本译照录；是否原拟写作 (I,VIII)，仅据现存见证不能唯一确定。总述 Lefschetz 公式的 (VII,X) 亦保留，第 VII 报告确有 Gauss--Bonnet 公式。

\hypertarget{collation-FRONT-02}{}

\paragraph*{FRONT-02。}

校勘注：法文导言第 VI 页误作 traces commutatives（交换迹）；据 III B 导言及第5节对同一报告 XI 和 Stallings 理论的说明，补正为非交换迹。

\hypertarget{collation-INDEX-01}{}

\paragraph*{INDEX-01。}

校勘注：法文原索引第 484 页此项作 VII 1.2；相应定义实际在第 VIII 报告第1、2节。下标 $D$ 表示其中的全子范畴。

\hypertarget{collation-XV-S1-new}{}

\paragraph*{XV-S1-new。}

校勘注：原文这一辅助断言未写出 $R$ 的平坦性；为使约化后的列仍然正合，应补此条件。在上面的递降构造中，$R_0=Z_0^{i+1}$ 为投射模，可由引理3(ii)取其投射提升 $R$，因而所需平坦性成立。

原文的辅助断言只要求：“$Q_0$ 投射，且 $R$ 为满足 $T(R)\simeq R_0$ 的 $A$-模”。该任意提升的表述不成立。例如取 $A=\Z/4\Z$、$A_0=\F_2$、$P_0=0$、$Q_0=R_0=\F_2$ 和 $R=\F_2$；投射提升 $Q=\Z/4\Z$ 到 $R$ 的满射之核约化后非零，不可能恢复原来的 $P_0=0$。正文只使用实际递降构造中 $Q_0,R_0$ 都投射、且可取 $R$ 投射的情形，并写明分裂保证约化正合的理由。

\hypertarget{collation-XV-S2-new}{}

\paragraph*{XV-S2-new。}

校勘注：法文原书第475页此处的 $K\prime$ 漏下标 $\nu$，按本段逐层构造补出。
